\section{Theoretische Grundlagen}
\label{sec:grundlagen}

\subsection{Cross-Site Scripting}

XSS entsteht, wenn nicht vertrauenswürdige Daten eine Ausgabesenke
erreichen, die den Wert als HTML oder JavaScript interpretiert. Der
eingeschleuste Code läuft anschließend im Origin der legitimen
Anwendung und unterliegt damit nicht der Same-Origin-Sperre gegenüber
dieser Anwendung. Er kann sichtbare Inhalte lesen, die Oberfläche
verändern und Requests über den vorhandenen Anwendungskontext
auslösen. OWASP empfiehlt deshalb kontextgerechtes Output Encoding
und sichere Framework-Ausgabe als primäre Schutzmaßnahme~\cite{owaspXss}.

Die klassischen Varianten unterscheiden sich durch Herkunft und
Lebensdauer des Payloads:

\begin{itemize}
  \item \textbf{Stored XSS}: Der Payload wird persistent gespeichert
    und bei späteren Aufrufen erneut ausgeliefert.
  \item \textbf{Reflected XSS}: Der Payload stammt aus dem aktuellen
    Request und wird unmittelbar zurückgegeben.
  \item \textbf{DOM-based XSS}: Die unsichere Verarbeitung erfolgt im
    Browser, beispielsweise über \texttt{innerHTML}.
\end{itemize}

Der im Scheduler gefundene Fall kombiniert Stored und DOM-based XSS:
Der Wert liegt dauerhaft in der Datenbank, wird über die API geladen
und erst beim Rendern durch \texttt{vue-cal} als HTML in das DOM geschrieben.

\subsection{Quellen, Senken und sichere Ausgabe}

Eine \emph{Quelle} ist ein kontrollierbarer Datenursprung,
beispielsweise ein Veranstaltungsname, Gruppenname oder
Reservierungslabel. Eine \emph{Senke} ist eine Operation, die Daten
als aktiven Inhalt interpretiert. Typische Senken sind
\texttt{innerHTML}, Vue-\texttt{v-html}, \texttt{insertAdjacentHTML}
oder ein Bibliotheksrenderer mit vergleichbarem Verhalten.

Vue-Interpolation mit doppelten geschweiften Klammern escaped HTML
standardmäßig und erzeugt Text statt aktiver
DOM-Strukturen~\cite{vueSecurity}. Genau diese Eigenschaft wird beim
Scheduler-Fix genutzt:

\begin{lstlisting}[caption={Sichere Ausgabe eines Kalenderereignisses}]
<template v-slot:event="{ event }">
    <div class="vuecal__event-content">
        <strong>{{ event.title }}</strong>
        <div>{{ event.content }}</div>
    </div>
</template>
\end{lstlisting}

Eine Content Security Policy (CSP) ist eine zusätzliche
Schutzschicht, ersetzt aber nicht die Behebung der unsicheren
Senke~\cite{owaspCsp}. Serverseitige Validierung verbessert
Datenqualität und begrenzt Payloads, kann kontextgerechtes Escaping
ebenfalls nicht ersetzen.

\subsection{Same-Origin Policy und CORS}

Eine Origin besteht aus Schema, Host und Port. Deshalb sind
\texttt{http://localhost:8078} und \texttt{http://localhost:8079}
verschiedene Origins. Die Same-Origin Policy verhindert
grundsätzlich, dass JavaScript einer Origin Antworten einer fremden
Origin liest. Ein Request kann dabei teilweise trotzdem gesendet
werden; blockiert wird vor allem der Zugriff des Skripts auf die Antwort.

CORS ist die kontrollierte Lockerung dieser Browserregel. Der Server
liefert dafür unter anderem folgende Header~\cite{mdnCors}:

\begin{itemize}
  \item \texttt{Access-Control-Allow-Origin}: erlaubte lesende Origin,
  \item \texttt{Access-Control-Allow-Credentials}: Freigabe
    browserseitiger Credentials,
  \item \texttt{Access-Control-Allow-Methods}: erlaubte HTTP-Methoden,
  \item \texttt{Access-Control-Allow-Headers}: erlaubte Request-Header,
  \item \texttt{Access-Control-Expose-Headers}: für JavaScript
    lesbare Response-Header.
\end{itemize}

Requests mit beispielsweise \texttt{Authorization} oder
\texttt{Content-Type: application/json} lösen üblicherweise einen
Preflight-Request mit \texttt{OPTIONS} aus. CORS muss deshalb vor der
Spring-Security-Authentifizierung verarbeitet werden, da der
Preflight noch keine Benutzer-Credentials enthält~\cite{springCors}.

\subsection{CORS ist keine Authentifizierung}

CORS vergibt keine fachlichen Rechte und schützt eine API nicht vor
direkten Requests außerhalb des Browsers. Authentifizierung
beantwortet, wer ein Nutzer ist; Autorisierung beantwortet, was
dieser Nutzer tun darf. CORS beantwortet lediglich, ob fremdes
Browser-JavaScript eine Antwort lesen darf.

Ebenso wichtig ist die Abgrenzung beim Scheduler: Das Frontend setzt
den Bearer-Token explizit im \texttt{Authorization}-Header. Eine
fremde Webseite erhält diesen Token nicht automatisch allein durch
eine offene CORS-Konfiguration. Ein Token-Leak setzt eine zusätzliche
Schwachstelle voraus, etwa XSS, eine unsichere Token-Antwort oder
falsch eingesetzte Cookie-Authentifizierung.
